  a=3^k J(s),\qquad k\ge0.
\]
We call $s$ its \emph{coefficient source}.  Every positive state has unique
canonical coordinates
\[
 q=Q_r^e(3^k J(s)).
\]
Indeed, if $q$ is even, take $e=0$, $r=v_2(q)$, and $a=q/2^r$; if $q$ is
odd, take $e=1$, $r=v_2(q+1)$, and $a=(q+1)/2^r$.  In both cases $a$ is
positive and odd, and its unique factorization above determines $k$ and
$s$.  We write $\rho(q)=s$ for this state source.

\begin{lemma}[Ordinary side relation]\label{lem:ordinary-side}
Let $q>0$.  If
\[
 q\not\equiv1,3,12,14\pmod {16},
\]
then
\[
 B(A(q))\in\{A(B(q)),B(B(q))\}.
\]
If $q$ is in one of the four exceptional classes, then $A(q)$ is not in an
exceptional class.
\end{lemma}

\begin{proof}
Put $x=A(q)$, $r=R(x)=B(q)$, and $k=\asl(x)$.  In the
nonexceptional classes one has $r>0$: if $r=0$, then $x$ is wholly
alternating, and the two short words are $10_2,101_2$ while every longer
word ends in $0101$ or $1010$.  Directly, the short words come from
$q=1,3$, and
\[
\begin{aligned}
 A(q)\equiv10\pmod {16}&\Longrightarrow q\equiv1,12\pmod {16},\\
 A(q)\equiv5\pmod {16}&\Longrightarrow q\equiv3,14\pmod {16}.
\end{aligned}
\]
Thus $r=0$ places $q$ in an exceptional class.  A direct residue check gives
\[
\begin{array}{c|c}
 k& q\pmod {16}\\ \hline
 1&0,2,5,7,8,10,13,15\\
 2&4,6,9,11.
\end{array}
\]
Thus every nonexceptional residue has $k\in\{1,2\}$.  We now treat these
two actual suffix lengths, so no fixed modulus is being used to decide an
unbounded suffix.

When $k=1$, the deleted bit equals the last bit
$\epsilon=r\bmod2$, and $x=2r+\epsilon$.  Directly,
\[
 A(x)=2A(r)+\epsilon.
\]
Thus $A(x)$ is obtained from $A(r)$ by appending the bit $\epsilon$.
If this bit equals the last bit of $A(r)$, the maximal alternating suffix
stops at the boundary and $R(A(x))=A(r)$; otherwise it continues through
exactly the maximal alternating suffix of $A(r)$ and
$R(A(x))=R(A(r))=B(r)$.

When $k=2$, the deleted word is $01$ if $r$ is even and $10$ if $r$ is
odd.  Direct substitution gives
\[
 A(x)=
 \begin{cases}
  4A(r)+2,&r\text{ even},\\
  4A(r)+1,&r\text{ odd}.
 \end{cases}
\]
Hence $A(x)$ is obtained from $A(r)$ by appending respectively $10$ or
$01$.  The same boundary test leaves either $A(r)$ or
$R(A(r))=B(r)$.  Since $R(A(x))=B(A(q))$, the first assertion follows.

For completeness, if $q=16u+c$ with
$c\in\{1,3,12,14\}$, then
\[
 A(q)\bmod16\in
 \begin{cases}
  \{2,10\},&c\in\{1,12\},\\
  \{5,13\},&c\in\{3,14\}.
 \end{cases}
\]
None of these residues is exceptional.
\end{proof}

\begin{lemma}[Source-boundary selector]\label{lem:source-selector}
For $s>0$, define
\[
 \alpha(s)=1-\bigl(\lfloor s/2\rfloor\bmod2\bigr).
\]
Factor
\[
 3J(s)+1-2e=2^v J(t),\qquad e\in\{0,1\}.
\]
Then
\[
 (v,t)=
 \begin{cases}
  (1,A(s)),&e=\alpha(s),\\
  (v,B(s))\text{ with }v\ge2,&e=1-\alpha(s).
 \end{cases}
\]
\end{lemma}

\begin{proof}
Put $x=A(s)$.  Its parity is $1-\alpha(s)$.  Since $J(s)=2x+1$,
\[
 3J(s)+1-2e=2(3x+2-e).
\]
For $e=\alpha(s)$ the quantity in parentheses is odd and equals $J(x)$:
it is $3x+1$ when $x$ is even and $3x+2$ when $x$ is odd.  This gives
$v=1$ and $t=A(s)$.

For the opposite phase $e=x\bmod2$, the quantity in parentheses is
\[
 3x+2-(x\bmod2)=3x+1+\bigl(1-(x\bmod2)\bigr).
\]
Lemma~\ref{lem:alt-factor} factors it as a positive power of two times
$J(R(x))=J(B(s))$.  Including the leading factor two gives $v\ge2$ and
$t=B(s)$.
\end{proof}

\begin{lemma}[Exceptional child-pair identity]\label{lem:exceptional-pair}
Let $q\equiv1,3,12,14\pmod {16}$ and put $L=B(q)$.  Then for some
$m\ge1$ and $\delta\in\{0,1\}$, the two children of $A(q)$ are exactly
\[
 \{Q_m^\delta(J(L)),Q_{m+1}^\delta(J(L))\}.
\]
\end{lemma}

\begin{proof}
Write $q=16u+c$.  The four direct formulas are
\[
\begin{array}{c|c|c|c}
 c&L&B(A(q))&A^2(q)\\ \hline
 1&R(6u)&18u+1&36u+3\\
 3&R(6u+1)&18u+4&36u+8\\
 12&R(6u+4)&18u+13&36u+27\\
 14&R(6u+5)&18u+16&36u+32.
\end{array}
\]
For $q=1$, one has $L=0$ and the two children of $A(q)=2$ are
$1=Q_1^1(J(0))$ and $3=Q_2^1(J(0))$, so the assertion is direct.
Otherwise put $z=6u,6u+1,6u+4,6u+5$ in the four rows.  Then $z>0$,
$L=R(z)$, and $B(A(q))=3z+1$.  Lemma~\ref{lem:alt-factor} gives
$3z+1=Q_m^\delta(J(L))$ for some $m\ge1$, where
$\delta=1-(z\bmod2)$.  The last column of the table is
$2B(A(q))+\delta$, hence is $Q_{m+1}^\delta(J(L))$.
\end{proof}

\begin{lemma}[Universal source bound]\label{lem:source-bound}
For every positive state $q$,
\[
  \rho(q)\le\frac{q-1}{6}.
\]
\end{lemma}

\begin{proof}
If $q=Q_r^e(3^k J(s))$, then $r\ge1$ and $e\le1$ give
$q\ge2J(s)-1$.  Since $J(s)\ge3s+1$, we have $q\ge6s+1$.
\end{proof}

\subsection{Two permanent negative regressions}

The first invalid shortcut in earlier manuscripts was
\[
 \text{smaller canonical coefficient}
 \quad\Longrightarrow\quad
 \text{smaller coefficient source}.
\]
It is false in the presence of powers of three.

\begin{proposition}[Coefficient/source separation]\label{prop:coeff-source-counter}
Let $s=1$, so $J(s)=5$, and put $a=3J(s)=15$, $e=1$.  The common boundary
child is
\[
 p=R(3a-e)=R(44)=22.
\]
Its canonical odd coefficient is $11<15$, but $11=J(3)$.  Hence its source
has increased from $1$ to $3$.
\end{proposition}

\begin{proof}
The displayed arithmetic is direct.  In binary, $44=101100_2$ and deleting
its maximal alternating suffix gives $10110_2=22$.  Since
$22=Q_1^0(11)$ and $11=J(3)$, the source is $3$.
\end{proof}

The second invalid shortcut compared a return after several expanding moves
with the numerical source before the streak.

\begin{proposition}[Pre-streak anchor separation]\label{prop:prestreak}
In the conjugated game,
\[
  20\xrightarrow{A}30\xrightarrow{A}45\xrightarrow{A}68,
  \qquad B(68)=25>20.
\]
Thus a contracting return after a multi-$A$ streak need not lie below the
source retained before the streak.
\end{proposition}

\begin{proof}
Each value follows from the definitions of $A$ and $B$.
\end{proof}

Neither implication is used later in this manuscript.

\section{Finite outcomes, witnesses, and retained tokens}

Every finite \WIN\ or \LOSS\ state has a canonical proof height.  Put
$h(0)=0$ and define
\[
 h(x)=1+\min\{h(y):y\text{ is a \LOSS\ child of }x\}
 \quad(x\text{ \WIN}),
\]
\[
 h(x)=1+\max\{h(y):y\text{ is a child of }x\}
 \quad(x\text{ \LOSS}).
\]
A finite state carried only to justify a local parent/child incidence is a
\emph{routing witness}.  A \emph{retained proof token} is a finite state
whose relation to the existing rank has also been proved: it either uses an
explicit one-shot initialization or replaces an already retained token by
proper proof-tree descendants.  Merely discovering that a state is finite
does not make it a retained token.

\begin{lemma}[Multiset token rank]\label{lem:token-rank}
For a finite multiset $M$ of proof heights, let
\[
  \Phi(M)=\mathop{\#}_{h\in M}\omega^h,
\]
where $\#$ is natural ordinal sum.  Replacing one occurrence of height $H$
by finitely many certified descendants of heights strictly below $H$
strictly decreases $\Phi$.
\end{lemma}

\begin{proof}
In Cantor normal form, the coefficient of $\omega^H$ decreases by one.
Every inserted monomial has smaller exponent, so no finite collection of
inserted terms can compensate for that loss.
\end{proof}

\begin{remark}[Why provenance is not cosmetic]
A finite \WIN\ state may have another \WIN\ child whose proof height is
not controlled by the canonical \LOSS\ branch used to prove the parent
\WIN.  Therefore a later finite state cannot be substituted into an
ordinal rank unless it is either initialized by a declared seed or proved to
be a proper descendant of a retained token.  This is the central audit rule
for the remaining sections.
\end{remark}

\section{A complete reverse fixed-tail progress lemma}

The following strengthens the fixed-tail statement used in Versions 5 and
6.  The old statement produced a finite witness but did not prove the
progress needed after that witness appeared.

\begin{lemma}[Reverse fixed-tail progress]\label{lem:fixed-tail-progress}
Let $c$ be positive and odd, $e\in\{0,1\}$, and $r\ge2$.  Put
\[
 X=Q_r^e(3c),\qquad
 G=Q_r^e(c),\qquad
 H=Q_{r+1}^e(c),\qquad
 Y=A(G).
\]
Then
\[
 \Moves(H)=\{X,Y\}.
\]
Let $w=B(X)$.  More precisely,
\[
 w=\begin{cases}
     B(Y),&r=2,\\
     A(Y),&r\ge3.
   \end{cases}
\]
If $X$ is \DRAW, exactly one of the following progress alternatives is
available:
\begin{enumerate}[label=(\roman*)]
\item $H$ is \DRAW; replacing $X$ by $H$ removes one factor of three from
      the displayed coefficient;
\item $H$ is \WIN, $Y$ is \LOSS, $w$ is a \WIN\ child of $Y$ with
      $h(w)\le h(Y)-1$, and the other child $A(X)$ is \DRAW.
\end{enumerate}
In alternative (ii), the continuation is therefore accompanied by the
strict token replacement $Y\mapsto w$ whenever $Y$ has legitimately been
initialized as a token.
\end{lemma}

\begin{proof}
Since $r+1\ge3$, Lemma~\ref{lem:long-tail} gives
\[
 A(H)=Q_r^e(3c)=X.
\]
For $r\ge3$ it also gives
$B(H)=Q_{r-1}^e(3c)=A(G)=Y$; for $r=2$ the shared boundary identity gives
the same equality.  Hence $\Moves(H)=\{X,Y\}$.

For $r=2$, $X=Q_2^e(3c)$ and $Y=Q_1^e(3c)$, so the shared boundary identity
gives $B(X)=B(Y)$.  If $r\ge3$, direct use of the long-tail formulas,
including the elementary $Q_2\mapsto Q_1$ row when $r=3$, gives
\[
 B(X)=Q_{r-2}^e(9c)=A(Y).
\]
This proves the common-child statement.

Assume $X$ is \DRAW.  Since $H$ has a \DRAW\ child, $H$ is not
\LOSS.  If $H$ is \DRAW, alternative (i) holds.  If $H$ is \WIN, one
of its children must be \LOSS.  Since $X$ is \DRAW, that child is
necessarily $Y$, so $Y$ is \LOSS.
The state $w$ is an actual child of this \LOSS\ state and is therefore
\WIN, with $h(w)\le h(Y)-1$.  It is also the $B$-child of $X$.  A \DRAW\
state has no \LOSS\ child and at least one \DRAW\ child.  Since $w$ is
\WIN, the other child $A(X)$ must be \DRAW.  This is alternative (ii).
\end{proof}

\begin{corollary}[Finite reverse factor countdown]\label{cor:reverse-countdown}
Starting from a displayed \DRAW\ state $Q_r^e(3^k a)$ with $r\ge2$ and
$k>0$, repeated use of Lemma~\ref{lem:fixed-tail-progress} either reaches a
factor-free displayed coefficient after finitely many reverse steps, or
